% --- CORRECCIÓN IMPORTANTE EN LA PRIMERA LÍNEA ---
% Agregamos "xcolor=table" para que funcione \rowcolor sin errores
\documentclass[aspectratio=169, 10pt, xcolor=table]{beamer}

% --- Configuración del Tema (estilo profesional como el ejemplo) ---
\usetheme{Madrid}
\usecolortheme{dolphin}

% --- Paquetes necesarios ---
\usepackage[utf8]{inputenc}
\usepackage[spanish]{babel}
\usepackage{amsmath, amssymb}
\usepackage{booktabs}
\usepackage{graphicx}
\usepackage{listings}
\usepackage{tikz}
\usetikzlibrary{arrows.meta, positioning, shapes.geometric}

% --- Ruta de Imágenes (crea la carpeta "imagenes" y coloca ahí tus archivos) ---
\graphicspath{{./imagenes/}}

% --- Estilo para código Python ---
\definecolor{codegreen}{rgb}{0,0.6,0}
\definecolor{codegray}{rgb}{0.5,0.5,0.5}
\definecolor{codepurple}{rgb}{0.58,0,0.82}
\definecolor{backcolour}{rgb}{0.95,0.95,0.92}
\lstset{
    language=Python,
    backgroundcolor=\color{backcolour},
    commentstyle=\color{codegreen},
    keywordstyle=\color{blue},
    numberstyle=\tiny\color{codegray},
    stringstyle=\color{codepurple},
    basicstyle=\ttfamily\scriptsize,
    breaklines=true,
    frame=single,
    showstringspaces=false
}

% --- Pie de página personalizado ---
\makeatletter

\setbeamertemplate{footline}{
  \leavevmode%
  \hbox{%
  \begin{beamercolorbox}[wd=.333333\paperwidth,ht=2.25ex,dp=1ex,center]{author in head/foot}%
    \usebeamerfont{author in head/foot}\insertshortauthor
  \end{beamercolorbox}%
  \begin{beamercolorbox}[wd=.333333\paperwidth,ht=2.25ex,dp=1ex,center]{title in head/foot}%
    \usebeamerfont{title in head/foot}Sesión 13
  \end{beamercolorbox}%
  \begin{beamercolorbox}[wd=.333333\paperwidth,ht=2.25ex,dp=1ex,right]{date in head/foot}%
    \usebeamerfont{date in head/foot}NLP \hfill \insertframenumber/\inserttotalframenumber \hspace*{0.3cm}
  \end{beamercolorbox}}%
  \vskip0pt%
}

\makeatother

% --- Datos de la portada ---
\title[SESIÓN 13 NLP]{\textbf{Sesión 13}}
\subtitle{Prompt Engineering y Optimización de Modelos}
\author[Prof. C. Sánchez]{Carlos César Sánchez Coronel}
\institute{\textbf{NLP}}
\date{}

% Logo opcional (descomentar si tienes la imagen)
% \titlegraphic{\includegraphics[height=1.5cm]{logo.png}}

% --- Eliminar la barra de navegación (opcional) ---
\setbeamertemplate{navigation symbols}{}

% --- Índice al inicio de cada sección ---
\AtBeginSection[]{
  \begin{frame}
    \frametitle{Contenido}
    \tableofcontents[currentsection]
  \end{frame}
}

% ============================================================================
% INICIO DEL DOCUMENTO
% ============================================================================
\begin{document}

% --- Portada ---
\begin{frame}
    \titlepage
\end{frame}

% --- Índice general ---
\begin{frame}{Contenido}
    \tableofcontents
\end{frame}

% ============================================================================
% SECCIÓN 1: PROMPT ENGINEERING E IN-CONTEXT LEARNING
% ============================================================================
\section{Prompt Engineering e In-Context Learning}

\begin{frame}{¿Qué es In-Context Learning (ICL)?}
    \begin{itemize}
        \item Capacidad de los LLMs para aprender a resolver tareas observando ejemplos en el prompt, sin actualizar sus pesos.
        \item Se diferencia del fine-tuning en que no requiere entrenamiento ni GPU.
        \item Estrategias principales: Zero-shot, Few-shot, Chain-of-Thought.
    \end{itemize}
    % Basado en Pregunta 1
\end{frame}

\begin{frame}{Zero-shot y Few-shot}
    \begin{columns}[t]
        \column{0.5\textwidth}
        \textbf{Zero-shot}
        \begin{itemize}
            \item Instrucción sin ejemplos.
            \item Ej: ``Clasifica el mensaje: ...''
            \item Depende del conocimiento pre-entrenado.
        \end{itemize}
        \column{0.5\textwidth}
        \textbf{Few-shot}
        \begin{itemize}
            \item Se proporcionan ejemplos en el prompt.
            \item Ej: ``Mensaje: ... $\rightarrow$ reclamo. Mensaje: ... $\rightarrow$ consulta. Ahora clasifica: ...''
            \item Guía el formato y estilo.
        \end{itemize}
    \end{columns}
    % Basado en Pregunta 1
\end{frame}

\begin{frame}{Chain-of-Thought (CoT)}
    \begin{block}{Idea}
        Pedir al modelo que ``piense paso a paso'' antes de dar la respuesta.
    \end{block}
    \begin{exampleblock}{Ejemplo}
        ``Problema: Si un tren viaja a 120 km/h durante 2.5 horas y luego a 80 km/h durante 1.5 horas, ¿qué distancia total recorre? Vamos a pensarlo paso a paso. Primero, calculamos la distancia del primer tramo: 120 km/h × 2.5 h = 300 km. Luego, la distancia del segundo tramo: 80 km/h × 1.5 h = 120 km. Finalmente, sumamos: 300 km + 120 km = 420 km. Respuesta: 420 km.''
    \end{exampleblock}
    \begin{itemize}
        \item Mejora el rendimiento en tareas de razonamiento lógico y matemático.
        \item Descompone problemas complejos en subtareas.
    \end{itemize}
    % Basado en Preguntas 1 y 2
\end{frame}

\begin{frame}{Self-Consistency en CoT}
    \begin{itemize}
        \item Generar múltiples cadenas de razonamiento independientes (con temperatura $>0$).
        \item Votar la respuesta más frecuente.
        \item Mejora la precisión a costa de mayor latencia y costo.
    \end{itemize}
    % Basado en Pregunta 15
\end{frame}

% ============================================================================
% SECCIÓN 2: FINE-TUNING EFICIENTE: LORA
% ============================================================================
\section{Fine-Tuning Eficiente: LoRA}

\begin{frame}{El problema del fine-tuning completo}
    \begin{itemize}
        \item Modelos con miles de millones de parámetros (7B, 13B, 70B).
        \item Fine-tuning completo requiere GPUs de alta memoria y tiempo.
        \item Riesgo de ``olvido catastrófico''.
    \end{itemize}
\end{frame}

\begin{frame}{LoRA (Low-Rank Adaptation)}
    \begin{block}{Idea}
        Los cambios en los pesos durante la adaptación tienen un ``rango intrínseco bajo''. En lugar de modificar la matriz original $W_0 \in \mathbb{R}^{d \times k}$, se añaden matrices de bajo rango entrenables.
    \end{block}
    \begin{columns}[t]
        \column{0.5\textwidth}
        \textbf{Arquitectura}
        \[
        W = W_0 + \Delta W = W_0 + B \cdot A
        \]
        donde $B \in \mathbb{R}^{d \times r}$, $A \in \mathbb{R}^{r \times k}$, $r \ll \min(d,k)$.
        \column{0.5\textwidth}
        \textbf{Ventajas}
        \begin{itemize}
            \item Parámetros entrenables: $r(d+k)$ vs $d \times k$.
            \item Evita olvido catastrófico.
            \item Adaptadores intercambiables para múltiples tareas.
        \end{itemize}
    \end{columns}
    % Basado en Preguntas 4 y 5
    % IMAGEN: Diagrama de LoRA
\end{frame}

\begin{frame}{Elección de hiperparámetros en LoRA}
    \begin{itemize}
        \item \textbf{Rango $r$}: típicamente 4, 8, 16, 32. Mayor $r$ = más capacidad, más parámetros.
        \item \textbf{Target modules}: normalmente las proyecciones de atención (\texttt{q\_proj}, \texttt{v\_proj}, \texttt{k\_proj}, \texttt{o\_proj}).
        \item \textbf{lora\_alpha}: escalado (similar a learning rate).
        \item \textbf{lora\_dropout}: regularización.
    \end{itemize}
    % Basado en Pregunta 6
\end{frame}

% ============================================================================
% SECCIÓN 3: CUANTIZACIÓN
% ============================================================================
\section{Cuantización}

\begin{frame}{¿Qué es la cuantización?}
    \begin{itemize}
        \item Reducir la precisión numérica de los pesos del modelo (ej. de FP32 a INT8).
        \item Reduce el uso de memoria y acelera la inferencia.
        \item Puede introducir una pequeña pérdida de calidad.
    \end{itemize}
\end{frame}

\begin{frame}{Formatos de cuantización}
    \begin{table}
        \centering
        \begin{tabular}{l|c|c|c}
            \toprule
            \textbf{Formato} & \textbf{Bits} & \textbf{Reducción} & \textbf{Pérdida típica} \\
            \midrule
            FP32 & 32 & 1x & - \\
            FP16 / BF16 & 16 & 2x & Mínima \\
            INT8 & 8 & 4x & Ligera \\
            4-bit (NF4) & 4 & 8x & Moderada (aceptable) \\
            \bottomrule
        \end{tabular}
    \end{table}
    \begin{itemize}
        \item NF4 (NormalFloat4) optimizado para pesos de redes neuronales.
        \item QLoRA: combina cuantización 4-bit con LoRA.
    \end{itemize}
    % Basado en PDF y Preguntas 7, 8
\end{frame}

\begin{frame}{Cuantización post-entrenamiento vs QAT}
    \begin{columns}[t]
        \column{0.5\textwidth}
        \textbf{PTQ (Post-Training Quantization)}
        \begin{itemize}
            \item Se cuantiza un modelo ya entrenado.
            \item Rápido, pero puede perder precisión.
        \end{itemize}
        \column{0.5\textwidth}
        \textbf{QAT (Quantization-Aware Training)}
        \begin{itemize}
            \item Se simula la cuantización durante el entrenamiento.
            \item Mejor calidad, pero requiere re-entrenamiento.
        \end{itemize}
    \end{columns}
    % Basado en Pregunta 17
\end{frame}

% ============================================================================
% SECCIÓN 4: EXPORTACIÓN A ONNX
% ============================================================================
\section{Exportación a ONNX}

\begin{frame}{ONNX (Open Neural Network Exchange)}
    \begin{itemize}
        \item Formato abierto para representar modelos de ML.
        \item Permite interoperabilidad entre frameworks (PyTorch, TensorFlow, etc.).
        \item ONNX Runtime optimiza la inferencia en múltiples plataformas (CPU, GPU, móvil).
        \item Ventajas: fusión de operadores, cuantización, kernels especializados.
    \end{itemize}
    % Basado en PDF y Pregunta 9
\end{frame}

\begin{frame}[fragile]{Exportar un modelo de PyTorch a ONNX}
    \begin{lstlisting}
import torch
import onnx

# Modelo de ejemplo
model = ...  # modelo entrenado
dummy_input = torch.randn(1, 3, 224, 224)

torch.onnx.export(
    model,
    dummy_input,
    "model.onnx",
    input_names=["input"],
    output_names=["output"],
    dynamic_axes={"input": {0: "batch_size"}, "output": {0: "batch_size"}}
)

# Verificar
onnx_model = onnx.load("model.onnx")
onnx.checker.check_model(onnx_model)
    \end{lstlisting}
    % Basado en Pregunta 9
\end{frame}

% ============================================================================
% SECCIÓN 5: COMPARACIÓN DE ESTRATEGIAS
% ============================================================================
\section{Comparación de Estrategias}

\begin{frame}{Prompt Engineering vs LoRA vs Fine-Tuning}
    \begin{table}
        \centering
        \begin{tabular}{l|c|c|c}
            \toprule
            \textbf{Característica} & \textbf{Prompt Eng.} & \textbf{LoRA} & \textbf{Fine-Tuning} \\
            \midrule
            Costo computacional & Muy bajo & Bajo & Alto \\
            Necesidad de datos & Ninguno & Pocos & Muchos \\
            Calidad esperada & Variable & Alta & Muy alta \\
            Flexibilidad & Alta (rápido) & Media & Baja (un modelo) \\
            \bottomrule
        \end{tabular}
    \end{table}
    % Basado en Pregunta 10
\end{frame}

% ============================================================================
% SECCIÓN 6: IMPLEMENTACIÓN CONCEPTUAL EN PYTHON
% ============================================================================
\section{Implementación Conceptual en Python}

\begin{frame}[fragile]{LoRA con PEFT (Hugging Face)}
    \begin{lstlisting}
from transformers import AutoModelForCausalLM
from peft import LoraConfig, get_peft_model

model = AutoModelForCausalLM.from_pretrained("base-model")

lora_config = LoraConfig(
    r=8,
    lora_alpha=32,
    target_modules=["q_proj", "v_proj"],
    lora_dropout=0.05,
    bias="none",
    task_type="CAUSAL_LM"
)

lora_model = get_peft_model(model, lora_config)
lora_model.print_trainable_parameters()
    \end{lstlisting}
    % Basado en Pregunta 11
\end{frame}

\begin{frame}[fragile]{Cargar modelo cuantizado en 4-bit}
    \begin{lstlisting}
from transformers import AutoModelForCausalLM, BitsAndBytesConfig

quant_config = BitsAndBytesConfig(
    load_in_4bit=True,
    bnb_4bit_compute_dtype=torch.float16,
    bnb_4bit_quant_type="nf4"
)

model = AutoModelForCausalLM.from_pretrained(
    "modelo-13B",
    quantization_config=quant_config,
    device_map="auto"
)
    \end{lstlisting}
    % Basado en Pregunta 12
\end{frame}

% ============================================================================
% SECCIÓN 7: PREPARACIÓN PARA ENTREVISTAS
% ============================================================================
\section{Preparación para Entrevistas}

\begin{frame}{Preguntas típicas}
    \begin{itemize}
        \item ¿Qué es In-Context Learning y en qué se diferencia del fine-tuning?
        \item Explica Zero-shot, Few-shot y Chain-of-Thought.
        \item ¿Qué es LoRA y por qué es más eficiente que el fine-tuning completo?
        \item ¿Qué es la cuantización? Compara FP16, INT8 y 4-bit.
        \item ¿Qué es ONNX y para qué sirve?
    \end{itemize}
    % Basado en PDF
\end{frame}

\begin{frame}{Preguntas avanzadas}
    \begin{itemize}
        \item ¿Cómo evita LoRA el olvido catastrófico?
        \item ¿Qué ventaja tiene poder intercambiar adaptadores para diferentes tareas?
        \item ¿Cómo afecta el rango $r$ en LoRA a la capacidad de adaptación?
        \item ¿Qué es QLoRA y por qué es útil?
        \item ¿Qué estrategias de prompting usarías para reducir alucinaciones?
    \end{itemize}
    % Basado en Preguntas 5, 6, 7, 13
\end{frame}

% ============================================================================
% SECCIÓN 8: CONEXIÓN CON EL RESTO DEL CURSO
% ============================================================================


% ============================================================================
% SECCIÓN 9: RESUMEN
% ============================================================================
\section{Resumen}

\begin{frame}{Lo que hemos aprendido}
    \begin{exampleblock}{Conceptos clave}
        \begin{itemize}
            \item \textbf{Prompt Engineering}: Zero-shot, Few-shot, Chain-of-Thought, Self-Consistency.
            \item \textbf{LoRA}: adaptación de bajo rango, eficiente en parámetros, evita olvido.
            \item \textbf{Cuantización}: FP16, INT8, 4-bit (NF4), QLoRA. Trade-off tamaño/calidad.
            \item \textbf{ONNX}: interoperabilidad y optimización de inferencia.
            \item \textbf{Comparación}: cuándo usar cada técnica según datos, costo y calidad.
            \item \textbf{Implementación}: PEFT, bitsandbytes.
        \end{itemize}
    \end{exampleblock}
\end{frame}

% --- Diapositiva final ---
\begin{frame}
    \centering
    \Huge \textbf{¡Gracias!}

\end{frame}

\end{document}